\chapter{渊民的水声}

水声在地底库深处响起。

这很不合理。

法厄同正在裂开，地底库干冷得像一具被掏空的骨头，可那声音确实是水。它没有管道破裂时的喷涌感，也不像液体滴落。那是一种有节奏的低响，像海浪从很远的地方穿过岩层，一下一下推到耳边。

洛安跟着瑟恩穿过下行通道。

通道尽头是一片破裂水穹。

它原本应该是渊民族设施的一部分。透明穹顶覆盖着一个巨大的水腔，水腔内有螺旋状平台和悬浮的生物灯。现在穹顶裂开，大部分水已经流失，只剩几条悬在低重力异常里的水带，像被扯断的蓝色绸缎。

几个渊民族幸存者站在水带之间。

他们的身体比人类修长，皮肤呈蓝灰色，颈侧鳃片不断开合。没有足够水分让他们保持状态，他们看起来像被陆地慢慢杀死。可他们仍围成一个圆，双手贴在一只深色容器上。

容器没有固定形状。

它像一团被水包着的黑色光。

每当渊民族发出低声，容器表面就泛起一圈涟漪。那些涟漪没有向外扩散，反倒向内收缩，像声音被压进更深的地方。

查理听见那声音时，忽然想起地球。

画面里没有城市，没有灵接舱，也没有查理熟悉的街道。那更接近纪录片里的深海。蓝鲸在黑暗海水中缓慢移动，发出低到人类胸腔都会共鸣的歌声。

我不知道为什么会想到蓝鲸。

这里不是地球。

这里是法厄同，一颗正在变成小行星带的远古星球。

可那水声太像某种还没有出生的歌。

一个渊民族转头看向洛安和瑟恩。他的眼睛没有眼白，像两枚深色水珠。

“震翼族。”他说，“你迟到了。”

瑟恩脸色变了。

“通道断了。节肢族在往这里来。”

“所有族群都在往这里来。”渊民说，“现在才想起海，是因为陆地已经保不住了。”

瑟恩没有反驳。

洛安看着那只容器，问：“那是什么？”

渊民没有立刻回答。他先看了看洛安，又像是看见洛安身体里另一个频率，目光停顿了一瞬。

“你身上有后来者。”

洛安僵住。

我也僵住。

瑟恩挡到洛安前面：“现在不是说这个的时候。”

“正因为现在不是时候，所以它才会来。”渊民说。

他把手重新贴回容器。

“这是巨鲸留下的种子。不是血肉，也不是卵。它是记忆、频率和一条生物道路。巨鲸无法在地球完整存续，所以它把自己压低，压到未来海洋能承受的尺度。”

洛安听不懂全部。

我也只懂一半。

可是“地球”和“未来海洋”两个词已经足够让我发冷。

“为什么交给海？”洛安问。

渊民低头，看着几乎干涸的水腔。

“因为海不会急着成为王。”

远处传来节肢族的脚步声。

渊民族幸存者们同时开始发声。低频水声层层叠加，破裂水穹里仅剩的水带开始震动。那只黑色容器逐渐缩小，从拳头大小变成一枚暗蓝色核，表面有类似鲸影的纹路一闪而过。

瑟恩伸手接住它。

就在种子落入他掌心的瞬间，整座水穹发出一声像叹息一样的低响。

我在那声音里听见了遥远的海。

也听见了尚未出生的蓝鲸。
